\section{Conclusion}
\label{sec:conclusion}

We presented CityFlow, an expert-ensemble multi-agent framework for personalized urban route planning. By decomposing the planning problem across 11 specialized experts with scene-aware dispatch, parallel wave execution, and review-rework feedback loops, our system achieves 90\% pass rate with an average score of 74.3/100 across 100 diverse scenarios (grade distribution: S:12, A:24, B:23, C:30, D:10, F:1), and 7.4/10 on 5 core scenarios with 5/5 pass rate. Speed-run scenarios achieve 100\% pass rate (20/20), while destination scenarios remain the weakest type (68.3 avg, 5 D-grades). The ablation analysis identifies algorithmic time smoothing as the most critical component, confirming that LLM outputs require deterministic validation for physical constraints.

Three design principles emerge from our systematic optimization journey---from 65.2 to 81.2 points across 100 scenarios and seven rounds of targeted fixes, validated by the final 100-scenario evaluation (90\% pass, 74.3 avg):

\begin{enumerate}[leftmargin=*]
    \item \textbf{Prompt before architecture}: Six lines of refusal prompt replaced a dedicated feasibility-gate component, suggesting that prompt engineering should be the first resort, not the last.
    \item \textbf{LLM proposes, algorithm validates}: Temporal and spatial constraints require deterministic verification; LLM semantic understanding alone is insufficient for physical reasoning.
    \item \textbf{Evaluation methodology matters}: Two of our seven optimization rounds fixed measurement errors, not algorithm bugs---evaluation methodology errors can masquerade as algorithm failures.
\end{enumerate}

Future work includes: (1)~integrating real-time data sources (traffic APIs, restaurant availability); (2)~expanding to multi-city support with transferable expert configurations; (3)~exploring lightweight learned routers to replace rule-based expert dispatch; and (4)~conducting user studies to validate that CityFlow routes match human satisfaction criteria.
